\documentclass[10pt,a4paper]{amsart}
\usepackage[margin=19mm]{geometry}
\usepackage{amsmath,amssymb,mathtools}
\usepackage[scr=rsfs,cal=euler]{mathalfa}
\usepackage{xfrac}
\usepackage[unicode,hidelinks,bookmarks=false]{hyperref}
\newcommand{\ie}{i.e.\ }
\newcommand{\cf}{cf.\ }
\newcommand{\resp}{respectively\ }
\newcommand{\Subsection}{\subsection}
\providecommand{\nobreakdash}{\nobreak\hbox{-}}
\setlength{\emergencystretch}{2em}
\multlinegap0pt
\usepackage{bm}
\usepackage{bbm}
\usepackage{dsfont}
\usepackage{xspace}
\usepackage{cancel}
\usepackage{mathtools}
\usepackage{amsthm}
\makeatletter
\@ifundefined{lemma}{\newtheorem{lemma}{Lemma}}{}
\makeatother
\makeatletter
\expandafter\ifx\csname Assumptions\endcsname\relax
\newenvironment{Assumptions}% Definition of assumptions
\fi
\expandafter\ifx\csname Definitions\endcsname\relax
\newenvironment{Definitions}% Definition of definitions
\fi
\newcommand{\R}{\mathbb{R}} \newcommand{\N}{\mathbb{N}}
\newcommand{\pt}{\partial_t}
\newcommand{\px}{\partial_x }
\DeclareMathOperator{\tv}{TV}
\makeatother
\title[On the structure of optimal solutions of conservation laws a]{On the structure of optimal solutions of conservation laws at a junction with one incoming and one outgoing arc}
\author{Fabio Ancona, Annalisa Cesaroni, Giuseppe Maria Coclite, Mauro Garavello}
\date{Source: arXiv:2507.10090v1 (CC BY 4.0)}
\begin{document}
\maketitle

\noindent\emph{Source and licence.} On the structure of optimal solutions of conservation laws at a junction with one incoming and one outgoing arc. Version arXiv:2507.10090v1 (CC BY 4.0), distributed under the Creative Commons Attribution 4.0 International licence, as recorded in the arXiv licence metadata for this exact version and shown on the abstract page https://arxiv.org/abs/2507.10090v1. Full rights notice: licenses/arxiv-2507.10090-CC-BY-4.0.txt.\par\smallskip

\noindent\textit{Editorial scope.} Counted unit: the Lemma whose source label is \texttt{lemmafdecrescente} in the article (source0.tex lines 2764--2784, source label \texttt{lemmafdecrescente}), delivered together with the article's own complete original proof of it (source0.tex lines 2785--2898, one \texttt{proof} environment of 114 lines). No mathematics is written, repaired or re-derived: statement and proof are the article's own text line for line. Required context retained uncounted: eq:CP-classic (lines 355--359); xbar (lines 2758--2758). Every definition, display and label of those lines is retained unchanged. Omitted from this package and not redistributed: 1--354, 360--2757, 2759--2763, 2899--4645. Nothing inside a retained excerpt is omitted or altered: every retained body is delivered exactly as the article's lines are written, so no omission receipt of any kind accompanies this package. Citations: the retained excerpts carry no citation command at all, so no bibliography entry of the article is transcribed and no reference is redistributed. Reference adaptation: none was needed and none was made; every reference inside the delivered unit resolves inside it, so no unresolved reference or citation is produced. Printed slips kept literal: no printed word, symbol, hypothesis or number of the counted statement or of its proof was altered. Layout and class: the article's own class is \texttt{article} with packages including a4wide, showkeys, graphicx, hyperref, amsfonts, amsthm, amsmath, amssymb, mathtools, cleveref, soul, tikz, pgfplots, pgf; the wrapper is the lane's pinned \texttt{amsart} editorial wrapper at 10pt on A4, which restates the theorem-like environments the retained text needs and adds the article's own macro definitions that the retained text needs (\texttt{\textbackslash R}, \texttt{\textbackslash pt}, \texttt{\textbackslash px}, \texttt{\textbackslash tv}), with extra wrapper packages (bm, bbm, dsfont, xspace, cancel, mathtools). The delivered numbering is the unit's own. Assets: no retained span contains a figure environment, a graphics-inclusion command, a PDF-page inclusion, a table or a TikZ picture; no image file is copied into this package, the package declares no asset dependency and no image file is redistributed with it. Licence: version arXiv:2507.10090v1 of this article is distributed under the Creative Commons Attribution 4.0 International licence, as recorded in the arXiv licence metadata for this exact version and shown on the abstract page https://arxiv.org/abs/2507.10090v1. A full rights notice is delivered at licenses/arxiv-2507.10090-CC-BY-4.0.txt. Characters: every retained excerpt is pure ASCII, so the delivered engine, XeLaTeX with its default Latin Modern text font, reports no missing character.
\begin{equation}
  \label{eq:CP-classic}
\begin{cases} \pt u + \px f(u) = 0, \quad & x \in \R, \, t > 0,\\ 
    u(0,x)= \overline u(x), & x \in \R,\end{cases}
\end{equation}

\begin{equation}\label{xbar} \bar x:= \sup\left\{x \in \R: \overline u(x) \le \theta\right\}.\end{equation}

\begin{lemma}\label{lemmafdecrescente}
  Assume that $\overline u $ is monotone non-decreasing and let  $\widetilde u$ be the unique entropy solution 
  to~\eqref{eq:CP-classic}.
  Then the following holds:
\begin{enumerate} \item for every $\hat x \in \R$, the map $t \mapsto f(\widetilde{u}(t, \hat{x}))$ is
  non-increasing,
  \item $f(\widetilde{u}(T^-, 0))\geq f(\widetilde{u}(T, 0^+)), f(\widetilde{u}(T, 0^-))$.\end{enumerate}

In particular, denoting $\widetilde \gamma=f(\widetilde u(\cdot,0))$, we have that 
\begin{align*} \mathcal F_{[0, T]}(\widetilde\gamma)
     = &
      \tv_{(0, T)}(f(\widetilde u(\cdot, 0)))
    +
    |f(\overline u(0^-))- \widetilde\gamma(0^+)| + 
    |f(\overline u(0^+))- \widetilde \gamma(0^+)| \\
    &+ 2\widetilde\gamma(T^-)-f(\widetilde{u}(T, 0^+))-f(\widetilde{u}(T, 0^-))
\\ = &  \widetilde \gamma(0^+)- \widetilde \gamma(T^-) +
    |f(\overline u(0^-))- \widetilde\gamma(0^+)| + 
    |f(\overline u(0^+))- \widetilde \gamma(0^+)| \\
    &+ 2\widetilde\gamma(T^-)-f(\widetilde{u}(T, 0^+))-f(\widetilde{u}(T, 0^-)).\end{align*}  
\end{lemma}

\begin{proof}
We start proving that $t\to f(\widetilde u(t, \hat x))$ is non-increasing.   Without loss of generality we assume 
that $\hat{x} = 0$.

  By contradiction assume that there exist $0 < t_1 < t_2$ such that
  $f(\widetilde{u}(t_1, 0)) < f(\widetilde{u}(t_2, 0))$. We may assume without loss of generality that they are continuity points. 
  Note that the Lax condition implies that the map $t \mapsto f(\widetilde{u}(t, 0))$
  can have only downward jumps. Hence, we may assume that there exists  a continuity point $\bar t \in (t_1, t_2)$ such that 
      \begin{equation}
        \label{eq:intermediate-flux}
        f(\widetilde{u}(t_1, 0)) < f(\widetilde{u}(\bar t, 0)) < f(\widetilde{u}(t_2, 0)). 
      \end{equation}
 
  
  %the map 
  %$t \mapsto f(\widetilde{u}(t, 0))$ is continuous in $(t_1, t_2)$ and that both
 % $t_1$ and $t_2$ are points of %continuity for $\widetilde{u}(t, 0)$.
  Denote with $\xi_1, \xi_2$ the backwards generalized characteristics exiting from
  $(t_1, 0)$ and $(t_2, 0)$. We have some possibilities.
  \begin{itemize}
      \item[(i)] $\xi_1(0) < \xi_2(0)$. In this case we necessarily deduce that $\xi_2(0) > 0$, otherwise
      $\xi_1$ and $\xi_2$ cross each other. Thus, $\widetilde{u}(t_2, 0) > \theta$.
      The monotonicity of $\overline u$ implies that
      $\widetilde{u}(t_1, 0) = \overline u(\xi_1(0)) \le \overline u(\xi_2(0)) = \widetilde{u}(t_2, 0)$.
      Since $f(\widetilde{u}(t_1, 0)) < f(\widetilde{u}(t_2, 0))$ we deduce that
      $\widetilde{u}(t_1, 0) < \theta$ and $\xi_1(0) < 0$.

      
     % This is possible since $t \mapsto f(\widetilde{u}(t, 0))$ is continuous in $(t_1, t_2)$.
      Denote with $\bar \xi$ the backward generalized characteristic exiting from $(\bar t, 0)$.\\
      If $\bar \xi(0) \le 0$, then $\bar \xi(0) < \xi_1(0)$ (otherwise $\bar \xi$ crosses $\xi_1$)
      and so $\widetilde{u}(\bar t, 0) \le \widetilde{u}(t_1, 0)<\theta $, 
      violating~\eqref{eq:intermediate-flux}.\\
      If $\bar \xi(0) > 0$, then $\bar \xi(0) \le \xi_2(0)$ (otherwise $\bar \xi$ crosses $\xi_2$)
      and so $\widetilde{u}(t_1, 0) < \theta < \widetilde{u}(\bar t, 0) \le \widetilde{u}(t_2, 0)$, 
      violating~\eqref{eq:intermediate-flux}.

      Therefore, this case is not possible.
    
    \item[(ii)] $\xi_1(0) = \xi_2(0)$. So in  $\xi_1(0)$ the initial data $\overline u$ would have a downward jump, 
    which is not possible by the assumption on monotonicity. 
    %  $\widetilde{u}(t_1, 0) = \widetilde{u}(t_2, 0)$.

    \item[(iii)] $\xi_1(0) > \xi_2(0)$. In this case we necessarily deduce that $\xi_2(0) < 0$, otherwise
      $\xi_1$ and $\xi_2$ cross each other. Thus, $\widetilde{u}(t_2, 0) < \theta$.
      The monotonicity of $\overline u$ implies that
      $\widetilde{u}(t_1, 0) = \overline u(\xi_1(0)) \ge \overline u(\xi_2(0)) = \widetilde{u}(t_2, 0)$.
      Since $f(\widetilde{u}(t_1, 0)) < f(\widetilde{u}(t_2, 0))$ we deduce that
      $\widetilde{u}(t_1, 0) > \theta$ and $\xi_1(0) > 0$.

      Consider now $\bar t \in (t_1, t_2)$ such that~\eqref{eq:intermediate-flux} holds 
      and $(\bar t, 0)$ point of continuity for $\widetilde{u}$.
      % This is possible since $t \mapsto f(\widetilde{u}(t, 0))$ is continuous in $(t_1, t_2)$.
      Denote with $\bar \xi$ the backward generalized characteristic exiting from $(\bar t, 0)$.\\
      If $\bar \xi(0) \ge 0$, then $\bar \xi(0) > \xi_1(0)$ (otherwise $\bar \xi$ crosses $\xi_1$)
      and so $\widetilde{u}(\bar t, 0) \ge \widetilde{u}(t_1, 0)$, 
      violating~\eqref{eq:intermediate-flux}.\\
      If $\bar \xi(0) < 0$, then $\bar \xi(0) \ge \xi_2(0)$ (otherwise $\bar \xi$ crosses $\xi_2$)
      and so $\widetilde{u}(t_2, 0) \le \widetilde{u}(\bar t, 0) < \theta <  \widetilde{u}(t_1, 0)$, 
      violating~\eqref{eq:intermediate-flux}.

      Therefore, this case is not possible. 
  \end{itemize}
So, the proof of 1. is concluded.
\smallskip

Now we show that $f(\widetilde{u}(T^-, 0))\geq f(\widetilde{u}(T, 0^+)), f(\widetilde{u}(T, 0^-))$. We show that 
$f(\widetilde{u}(T^-, 0))\geq  f(\widetilde{u}(T, 0^-))$ (the case $f(\widetilde{u}(T^-, 0))\geq f(\widetilde{u}(T, 0^+))$ 
is completely analogous). 

Let us consider an increasing  sequence $x_n\to 0^-$   such that $(T,x_n)$  are continuity points for $\widetilde u$ 
and $f(\widetilde{u}(T, x_n))\to f(\widetilde{u}(T, 0^-))$. We denote with $\xi_n$ the backward genuine characteristic 
passing through $(T,x_n)$. 

 
We distinguish three cases.
\begin{itemize}
\item[(i)] There exists $\bar t<T$ such that all points $(t,0)$ for $t\in (\bar t, T)$ 
are discontinuity points for $\widetilde u$. In this case there exists a shock located in $x=0$ in the interval 
$(\bar t, T)$, so $\widetilde u(t, 0^-)\leq \theta$ and $\widetilde u(t, 0^+)\geq \theta$ for all $t\in  (\bar t, T)$. 
We consider a sequence $t_n\to T^-$ and the backward minimal characteristics $\widetilde \xi_n$ from $(t_n, 0)$.  
Then necessarily $\xi_n(0)\leq \widetilde \xi_n(0)\leq \bar x$, where $\bar x$ is defined in \eqref{xbar}. 
Then, we conclude by the monotonicity of $\overline u$, we get that 
$\widetilde u(T,x_n)\leq \widetilde u(t_n, 0^-)\leq \theta$, from which we deduce 
 $ f(\widetilde{u}(T, 0^-))\leq f(\widetilde{u}(T^-, 0))$.
\item[(ii)] There exists $\bar t<T$ such that  there exists a dense subset  of $ (\bar t, T)x\{0\}$ 
of continuity points for $\widetilde u$. 

Assume that  
$\widetilde u(t, 0)\leq \theta$ for $t\to T^-$ and $(t,0) $ continuity point. Let us denote $\widetilde\xi$ is the backward 
genuine characteristic passing through $(t,0)$.

In this case, $\widetilde u(T, x_n)\leq \theta$, and we proceed exactly as in case $1$. 
\item[(iii)] There exists $\bar t<T$ such that  there exists a dense subset  of $ (\bar t, T)\times\{0\}$ 
of continuity points for $\widetilde u$. 

Assume that  
$\widetilde u(t, 0)> \theta$ for $t\to T^-$ and $(t,0) $  continuity point. Let us denote $\widetilde\xi$ 
is the backward genuine characteristic passing through $(t,0)$, for $t\in (\bar t, T)$, such that $(t,0)$ continuity point.

If $\widetilde u(T, x_n)>  \theta$, then we proceed as in case $1$, by symmetry.

It remains to consider the case $\widetilde u(T, x_n)\leq \theta$. Assume that there exists $\tau\in (t, T)$ 
such that $\widetilde \xi(\tau)=0$. Then this would imply $f'(\widetilde(\tau, 0^-))<\widetilde\xi'(\tau)$, 
in contradiction with Lax conditions. Moreover, due to the fact that $\widetilde u(T, 0^-)\leq \theta$, 
we conclude that necessarily $\widetilde \xi(T)=0$ with $\widetilde \xi'(T)\geq 0$. This means that $\widetilde{u}$ 
admits a shock discontinuity
at $(T,0)$, connecting the left state $\widetilde{u}(T, 0^-)$ with the right state $\widetilde{u}(T^-, 0)$,
that has slope $\widetilde \xi'(T)\geq 0$,
which implies that $f(\widetilde{u}(T^-, 0))\geq  f(\widetilde{u}(T, 0^-))$.
\end{itemize}
 

\end{proof}

\end{document}
